\scanpage{32}
16)\quad\uline{Classes fondamentales globales}

On fixe un entier $n>0$, et on se
\deleted{Soit $f:X\longrightarrow Y$ morphisme}
borne à des faisceaux de $n$-torsion (i.e. Modules sur $A:\mathbb Z/n\mathbb Z$),
et des préschémas à car. résiduelles premières à $n$
i.e. tels que $n.1_{\mathcal O_X}$ est inversible dans $\mathcal O_X$.

Si $f:X\mathbin{-}Y$ est un morphisme lisse, on dira qu'il est
\inserted{purement} de dimension relative $d$ si ses fibres sont
\inserted{\unsure{pures|purement}} de dim $d$.
On \inserted{désigne} par $T^{(n)}_{X/Y}$ ou simplement $T_{X/Y}$
s'il n'y a pas de confusion possible, le faisceau
$\mu_{nX}^{\otimes d}$.
[Si $f$ est lisse de dim \inserted{\unsure{relative loc. variable}} quelconque,
alors on définit de façon évidente $T_{X/Y}$ en le prenant sur les divers morceaux].

Nous regardons pour la suite des morphismes
\uline{lisses compactifiables purement de dim} $d$
$f:X\longrightarrow Y$.
Pour un tel morphisme, on peut définir un
\uline{\deleted{\unsure{hom}} \inserted{hom.} canonique}
\mspage{31}
